% Generated draft: requires mathematical and source review.
\documentclass[11pt]{article}
\usepackage[margin=25mm]{geometry}
\usepackage{fontspec}
\setmainfont{Latin Modern Roman}
\usepackage{amsmath,amssymb,mathtools,unicode-math}
\setmathfont{Latin Modern Math}
\usepackage{enumitem,needspace,xurl,hyperref}
\setlength{\parindent}{0pt}
\setlength{\parskip}{5pt}
\setlength{\emergencystretch}{3em}
\newcommand{\N}{\mathbb{N}}
\newcommand{\Z}{\mathbb{Z}}
\newcommand{\Q}{\mathbb{Q}}
\newcommand{\R}{\mathbb{R}}
\newcommand{\C}{\mathbb{C}}
\newcommand{\F}{\mathbb{F}}
\newcommand{\A}{\mathbb{A}}
\newcommand{\PP}{\mathbb{P}}
\newcommand{\E}{\mathbb{E}}
\newcommand{\eps}{\varepsilon}
\newcommand{\dd}{\,\mathrm{d}}
\begin{document}
\section*{UnsolvedMath ID: 30005181}
\textbf{Source problem code:} catalogue entry
\subsection{Definitions and mathematical statement}
Under the usual smooth interpretation, an integral homology three-sphere is a closed, connected, oriented smooth three-manifold \(Y\) satisfying
\[
H_k(Y;\Z)\cong H_k(S^3;\Z)\qquad\text{for every }k\geq 0.
\]
Two such manifolds \(Y_0,Y_1\) are integrally homology cobordant if there is a compact, oriented smooth four-manifold \(W\) with
\[
\partial W=(-Y_0)\sqcup Y_1
\]
such that both boundary inclusions induce isomorphisms on homology with coefficients in \(\Z\) in every degree. Here \(-Y_0\) denotes reversal of orientation. The group \(\Theta^3_{\Z}\) consists of these homology-cobordism classes, with addition induced by oriented connected sum, identity \([S^3]\), and inverse \(-[Y]=[-Y]\).

The question is whether this group has a nonzero element of finite order:
\[
\exists [Y]\in\Theta^3_{\Z}\;\exists n\in\Z,\quad
n\geq 2,\qquad [Y]\neq [S^3],\qquad n[Y]=[S^3].
\]
Here \(n[Y]\) is the class of the connected sum of \(n\) copies of \(Y\).

The supplied record does not explicitly specify the smooth or topological category. The formulation above uses the usual smooth interpretation of this notation; confirmation of that convention from the cited report is required. The supplied label \emph{open} and the supplied assertion of source verification are curation metadata, not independent verification in this draft.
\subsection{Short English statement}
Does the integral homology-cobordism group of oriented integral homology three-spheres contain a nonzero element of finite order? Under the smooth interpretation, this asks whether some sphere has a nonzero class whose connected sum with itself finitely many times represents \([S^3]\). The category requires source confirmation.
\subsection{Sources}
\begin{itemize}
\item Supplied local UnsolvedMath registry record, permanent ID \(30005181\), source problem code \texttt{catalogue entry}, titled \emph{Torsion in the Integral Homology-Cobordism Group}.
\item Supplied report citation: \emph{Topologie (2023)}, Oberwolfach Reports, DOI 10.4171/owr/2022/34, pp.~1927--1986. Original Oberwolfach report: \url{https://doi.org/10.4171/owr/2022/34}.
\item The supplied original statement places the question in \emph{Homology cobordism and Heegaard Floer homology}, Jennifer Hom, joint work with Irving Dai, Kristen Hendricks, Matthew Stoffregen, Linh Truong, and Ian Zemke. No separate proposer is identified in the record.
\item Supplied literature source, labeled ``survey of the homology-cobordism group'': \url{https://arxiv.org/abs/2209.10735}.
\item Literature assessment dated exactly \texttt{2026-08-22}: the supplied curation labels the question open and states that existing work establishes large infinite-rank subgroups without settling the existence of torsion. This is supplied dated curation derived from the Oberwolfach record, not independently verified here; its time-sensitive status and cited claims require independent checking.
\item The supplied statement status is \texttt{corrected\_verified}, with the supplied qualification that the statement was reconstructed and checked against the official report. This draft does not independently confirm that verification. The citation says \emph{Topologie (2023)}, while the proposed year and workshop-year tag are 2022; these supplied dates are retained without reconciliation.
\end{itemize}
\end{document}
